%==============================================================================
% Anexos
%==============================================================================

\section{Diagrama de arquitectura del sistema}

El Anexo~A reune los diagramas de la arquitectura propuesta en notación UML. Incluye:

\begin{itemize}
  \item \textbf{Diagrama de despliegue}: los cinco componentes principales (N8N, FastAPI, PostgreSQL, Gemini~2.5 Flash, Twilio), el proxy inverso Nginx que, en el diseño, provee terminación TLS~1.2/1.3 en el perímetro del sistema (no incluido en el despliegue publicado), y las relaciones de comunicación entre ellos.
  \item \textbf{Diagrama de secuencia}: flujo extremo a extremo de un incidente desde su recepción en el canal de correo electrónico hasta la confirmacion al usuario, incluyendo las interacciones con el clasificador hibrido y la base de datos.
  \item \textbf{Diagrama de componentes}: organización interna del modulo Python por capas (routes, services, repositories, models, classifiers), con sus dependencias y relaciones.
\end{itemize}

Los diagramas se mantienen en formato editable (\texttt{.drawio}) en el repositorio del proyecto para facilitar su actualizacion futura. La versión en alta resolución se encuentra disponible en el directorio \texttt{docs/} del repositorio \parencite{Gomez2026_repositorio}.

\section{Repositorio de código fuente}

El código fuente del modulo Python, los scripts de migración de la base de datos PostgreSQL gestionados mediante Alembic, los archivos de configuración del despliegue en contenedores Docker y los flujos de integración continua sobre GitHub Actions se encuentran disponibles en el repositorio publico del proyecto \parencite{Gomez2026_repositorio}. El commit de referencia para la versión entregada al jurado se consigna mediante el tag \texttt{v1.0.0}. La creación de dicho tag sobre el commit definitivo entregado al jurado queda como acción pendiente al cierre de los changes c-76 y c-77; una vez creado, se registrará en este mismo apartado el hash completo del commit etiquetado como referencia verificable de la versión evaluada.

La estructura del repositorio sigue el patron estandar con directorios separados para el modulo de procesamiento (\texttt{Gestion\_Incidentes/}), la definición del flujo de orquestación (\texttt{Automatizacion\_Mesa\_de\_Ayuda.json}), las migraciones de base de datos (\texttt{alembic/}), las pruebas automatizadas (\texttt{tests/}), los artefactos de evaluación (\texttt{evaluation/}) y la documentación operativa (\texttt{docs/}). El archivo \texttt{README.md} incluye instrucciones detalladas para el despliegue local en menos de 15 minutos.

\section{Esquema de base de datos}

El esquema completo de la base de datos PostgreSQL se documenta mediante scripts SQL versionados que definen las siete entidades descritas en la Tabla~\ref{tab:entidades} del Capitulo~5, las restricciones de integridad referencial (claves foraneas con eliminacion en cascada condicional), los indices secundarios sobre los atributos de busqueda frecuente (fecha de creacion, sector responsable y estado actual del ticket), y la generacion de los identificadores de incidente. Las migraciones incrementales se gestionan mediante Alembic, lo que permite reproducir el estado de la base en cualquier punto histórico del desarrollo y revertir cambios cuando resulte necesario.

La migración inicial incluye la creacion de las tablas de catalogo (\texttt{sector}, \texttt{estado}, \texttt{canal\_origen}) con sus valores semilla, garantizando que un despliegue nuevo del sistema comience con los datos de referencia correctos. Las migraciones subsiguientes documentan la evolucion del esquema durante las seis iteraciones de desarrollo descritas en el Capitulo~4.

\section{Especificación OpenAPI de la interfaz REST}

La especificación completa de la interfaz REST se genera automáticamente mediante FastAPI conforme a la versión~3.1 de la especificación OpenAPI \parencite{Tiangolo2024_fastapi}. El documento resultante describe los puntos de entrada listados en la Tabla~5 del Capitulo~5, los esquemas de los cuerpos de solicitud y respuesta validados mediante Pydantic~v2, los códigos de estado posibles para cada método HTTP, y ejemplos representativos para cada caso de uso.

La interfaz se encuentra disponible adicionalmente en formato interactivo a través del componente Swagger UI integrado en el servidor, accesible bajo la ruta \texttt{/docs} durante la ejecución, y en formato estatico como archivo \texttt{openapi.json} en el directorio \texttt{docs/} del repositorio. La sincronización entre el código y la especificación estatica se verifica automáticamente en el pipeline de integración continua.

\section{Configuración del flujo N8N}

La exportación del flujo de orquestación en formato JSON incluye la totalidad de los 29 nodos configurados, sus parámetros operativos, las credenciales referenciadas mediante identificadores opacos sin exposición de valores, y los conectores entre nodos. El archivo \texttt{Automatizacion\_Mesa\_de\_Ayuda.json} en la raíz del repositorio contiene esta exportación. La importación de este artefacto en una instancia limpia de N8N~2.11.2 reproduce integramente el comportamiento operativo del sistema, sujeta únicamente a la provisión de las credenciales correspondientes a los servicios externos integrados (servidor IMAP de correo, cuenta Twilio y clave de API de Gemini) mediante el sistema de credenciales nativo de la plataforma.

\section{Corpus de validación}

El corpus de 200 incidentes utilizado para la validación experimental reside en el repositorio en formato JSON, en su forma pseudonimizada y versionada, bajo la ruta \texttt{data/corpus\_evaluacion\_pseudonimizado.json}. Cada caso incluye el identificador anonimo, la descripción cruda pseudonimizada, el canal de origen, el sector asignado por consenso humano, los sectores adicionales de verdad y el tiempo de registro medido en cada flujo. La categoría predicha por el clasificador y la confianza asociada no forman parte del corpus: se conservan en \texttt{evaluation/predicciones.json}, generado por el framework de evaluación. Se incluyen además las planillas de cronometraje del flujo manual y los registros automáticos generados por el flujo automatizado.

Por razones de privacidad y en cumplimiento de la Ley~25.326 \parencite{Ley25326_2000}, las descripciones textuales han sido pseudonimizadas previamente a su inclusion en el repositorio publico, reemplazando nombres propios, direcciones de correo y otros identificadores personales por etiquetas genericas.

\section{Documentación operativa}

La documentación operativa del sistema describe los procedimientos de despliegue inicial, configuración recurrente, respaldo y restauración, monitoreo activo y respuesta ante incidentes operativos. Se incluyen instrucciones detalladas para los entornos de desarrollo, preproduccion y produccion, así como una guia de resolución de incidencias frecuentes orientada al personal de operaciones de la organización adoptante. La documentación se mantiene actualizada en el repositorio del proyecto y se versiona conjuntamente con el código fuente, garantizando coherencia entre la versión desplegada y los procedimientos vigentes en cada momento del ciclo de vida del sistema.

\section{Trazabilidad de las correcciones del dictamen}

El registro de correcciones derivado del dictamen evaluador se conserva como artefacto versionado en el repositorio, bajo la ruta \texttt{docs/Tesis/Correcciones/dictamen-correcciones.md}. El documento compila, de forma acumulativa, cada corrección aplicada al manuscrito y las decisiones adoptadas sobre las observaciones que no admitían una resolución inmediata (por ejemplo, la convención estadística del estadístico \emph{W}, la reformulacion de la hipótesis principal y el tratamiento de los reintentos como limitación conocida). Se incluye aquí como ejemplo de trazabilidad: cada entrada del registro vincula una observación del dictamen con el archivo y la linea corregidos, de modo que la correspondencia entre la revisión solicitada y la revisión efectuada resulte auditable de extremo a extremo.

\newpage
